\documentclass[11pt]{article}
\usepackage[a4paper,margin=28mm]{geometry}
\usepackage{amsmath,amssymb,amsthm}
\usepackage[T1]{fontenc}
\usepackage{lmodern}
\usepackage{microtype}
\usepackage[hidelinks]{hyperref}
\usepackage{booktabs}
\newtheorem{theorem}{Theorem}
\newtheorem{proposition}{Proposition}
\title{An Improved Lower Bound for the Finite Zarankiewicz Number $z(13,23;3,3)$}
\author{Zijian Zeng\\
\small Institute of Computer Science and Digital Innovation, UCSI University\\
\small Kuala Lumpur, 56000, MALAYSIA\\
\small \texttt{zijianzeng@foxmail.com}; \texttt{1002266693@ucsiuniversity.edu.my}}
\date{}
\hypersetup{pdftitle={An Improved Lower Bound for the Finite Zarankiewicz Number z(13,23;3,3)},pdfauthor={Zijian Zeng}}

\begin{document}
\maketitle

\begin{abstract}
We give an explicit $13\times23$ binary matrix with 140 ones and no all-one
$3\times3$ submatrix. This improves the previously located lower bound 139.
For completeness, we also give a short independent proof of the known upper
bound 144. Consequently,
\[
140\le z(13,23;3,3)\le144.
\]
The exact value remains open.
\end{abstract}

\section{Introduction}
For positive integers $m,n,s,t$, let $z(m,n;s,t)$ be the maximum number of
ones in an $m\times n$ binary matrix containing no all-one $s\times t$
submatrix. Equivalently, it is the maximum number of edges in a bipartite
graph with parts of orders $m,n$ and no copy of $K_{s,t}$.

An exact result $z(13,22;3,3)=137$ was recently reported
\cite{afrasyab2026}; adding a column of weight two gives
$z(13,23;3,3)\ge139$. A public verification artifact gives the upper bound
$z(13,23;3,3)\le144$ and explicitly records the interval $139$--$144$
\cite{FiniteZarankiewiczClosures2026}. The general finite upper-bound context is developed by
Davies, Gill, and Horsley \cite{davies2026}.

The contribution here is the following one-edge improvement.

\begin{theorem}\label{thm:interval}
The finite Zarankiewicz number satisfies
\[
140\le z(13,23;3,3)\le144.
\]
\end{theorem}

The literature and exact-string searches used for the novelty check were
completed on 31 August 2026. Thus the priority statement is limited to the
sources indexed and cited here; the mathematical lower bound itself is
independently checkable from the displayed matrix.

\section{The 140-edge construction}
Read each of the following strings as one row of a binary matrix, from left to
right. The row weights are
$(11,11,11,11,11,11,12,11,11,11,11,12,6)$.
\begin{center}
\begin{verbatim}
11111111111000000000000
11111000000111110000100
11000111000111001110000
10100100110100101101100
10010010101010111011000
10001001011001010111100
01100001101010011100111
01010100011001101010011
01001010110100110110010
00110011010101011001001
00101110001011100101010
00011101100110000011111
01001100000000000100101
\end{verbatim}
\end{center}

\begin{proposition}\label{prop:witness}
The displayed matrix has 140 ones and contains no all-one $3\times3$
submatrix.
\end{proposition}

\begin{proof}
Summing the displayed row weights gives 140. For any three rows, direct
intersection of their supports has cardinality at most two. Indeed, exhaustive
inspection of all $\binom{13}{3}=286$ row triples gives maximum common support
two. This condition is equivalent to the absence of three columns on which
the same three rows are all one, and hence to the absence of an all-one
$3\times3$ submatrix.
\end{proof}

The first twelve rows are a relabeling-free copy of a 134-edge construction
distributed with \cite{FiniteZarankiewiczClosures2026}. The final row has support
$\{2,5,6,18,21,23\}$ in one-based indexing. Exhausting all $2^{23}$ possible
added-row supports over this fixed core shows that weight six is maximal for
this particular extension scheme. This fixed-core optimality is not a global
upper bound and is not needed for Proposition~\ref{prop:witness}.

\section{The upper bound}
We include the short argument underlying the upper endpoint in
Theorem~\ref{thm:interval}, independently of the software artifact
\cite{FiniteZarankiewiczClosures2026}.

Suppose that a feasible matrix has 145 ones. For column $j$, let $E_j$ be its
row support and $d_j=|E_j|$. Every row triple occurs in at most two columns, so
\begin{equation}\label{eq:capacity}
\sum_{j=1}^{23}\binom{d_j}{3}\le2\binom{13}{3}=572.
\end{equation}
For $0\le d\le13$,
\[
\binom d3\ge15d-70,
\]
with equality only for $d=6,7$. The total nonnegative penalty
$p(d)=\binom d3-15d+70$ is at most seven. The column count, degree sum 145,
and this penalty budget leave exactly the profiles
\[
6^{16}7^7,\qquad 5^1 6^{14}7^8,\qquad 6^{17}7^5 8^1.
\]
Their unused triple capacities
$s=572-\sum_j\binom{d_j}{3}$ are respectively $7,2,1$.

For a row $r$, put
\[
D_r=132-\sum_{j:r\in E_j}\binom{d_j-1}{2}\ge0.
\]
Double counting gives $\sum_rD_r=3s$. A row outside every exceptional
degree-five or degree-eight column receives only contributions 10 and 15, so
$D_r\equiv132\equiv2\pmod5$ and therefore $D_r\ge2$. In the three profiles,
at least $13,8,5$ rows are respectively clean. They force total deficits at
least $26,16,10$, whereas $3s$ equals $21,6,3$. Each case is impossible.
Thus 145 edges are impossible; deleting edges handles every larger weight and
proves the upper bound 144. Together with Proposition~\ref{prop:witness}, this
proves Theorem~\ref{thm:interval}.

\section{Verification and limitations}
The construction was checked in two independent exact ways: by all 286 row
triples and directly by all
$\binom{13}{3}\binom{23}{3}=506{,}506$ candidate submatrices. The degree-profile
classification in the upper-bound proof was also regenerated by an exact
integer enumerator. No floating-point computation or solver infeasibility
claim enters either endpoint.

Computational search was used to locate the final row and to implement
redundant checks. All theorem-bearing data and arguments are printed here or
use finite integer enumeration; responsibility for the claims and
presentation rests with the author.

The remaining gap is genuine: this work neither constructs 141--144 edges nor
certifies their impossibility. In particular, fixed-core extension optimality
must not be confused with the exact Zarankiewicz number.

\bibliographystyle{plain}
\bibliography{references}
\end{document}
